\documentclass[11pt]{article}
\usepackage{fontspec}
\setmainfont{Times New Roman}
\setsansfont{Arial}
\setmonofont{Consolas}
\usepackage{amsmath,amssymb,mathtools}
\usepackage{geometry}
\usepackage{microtype}
\newcommand{\mS}{\mathfrak{S}}
\newcommand{\mA}{\mathfrak{A}}
\newcommand{\mat}[1]{\begin{pmatrix}#1\end{pmatrix}}
\geometry{a4paper,margin=2.35cm}
\setlength{\parindent}{1.25em}
\setlength{\parskip}{0pt}
\makeatletter
\renewcommand\footnotesize{\@setfontsize\footnotesize{7.6}{8.6}\selectfont}
\makeatother
\emergencystretch=100em
\allowdisplaybreaks
\sloppy
\hbadness=10000
\vbadness=10000
\hfuzz=2pt
\vfuzz=10pt
\raggedbottom

\begin{document}

% Unit: Paper32 section 3 cyclic quaternion fields v001. Direct Ukrainian translation from the R122 German Paper32/Paper33 source slice, lines 482--516 / segments P32-S0023--P32-S0029.
\noindent Далі слід показати, що для нескінченно багатьох $n$ існують циклічні поля степеня $2^n$, які є мінімальними полями розщеплення кватерніонної групи.

\medskip

\noindent Нехай $p$ -- раціональне просте число, для якого при деякому $r>0$
\[
        (*)\qquad 2^r+1\equiv0\pmod p
\]
виконується; нехай $2^n$ буде найбільшою степенню числа $2$, що ділить $p-1$.

Ми спершу покажемо, що для нескінченно багатьох значень $n$ існують прості числа $p$. Для цього нехай $k$ -- довільне додатне ціле раціональне число, а $p$ -- простий дільник числа $2^{2^k}+1$. Тоді потрібна для $p$ конгруенція виконується при $r=2^k$. Крім того, $2\pmod p$ має порядок $2^{k+1}$, а тому для найбільшої степені $2^n$ числа $2$, що ділить $p-1$, число $n$ більше за $k$. Отже, у всякому разі існують як завгодно великі $n$, для яких є відповідні прості числа $p$.

Нехай $\varepsilon$ -- примітивний $p$-й корінь з одиниці. Тепер покажемо, що в $P(\varepsilon)$ число $-1$ подається як сума двох квадратів
\[
        -1=a^2+b^2=(a+ib)(a-ib)\qquad (a,b\text{ числа з }P(\varepsilon)).
\]
Це рівносильно тому, що в полі $(4p)$-х коренів з одиниці $P(\varepsilon,i)=P(\varepsilon\cdot i)$ існують числа, відносна норма яких щодо $P(\varepsilon)$ дорівнює саме $-1$.

Число $1+i\varepsilon^\nu$ $(\nu=1,2,\ldots,p-1)$ має відносну норму $(1+i\varepsilon^\nu)(1-i\varepsilon^\nu)=1+\varepsilon^{2\nu}$; отже, всі числа $1+\varepsilon^\nu$ $(\nu=1,2,\ldots,p-1)$ є відносними нормами чисел з $P(\varepsilon i)$ щодо $P(\varepsilon)$ як основного поля. Тому те саме справджується і для
\[
\Pi=\prod_{\nu=0}^{r-1}(1+\varepsilon^{2^\nu})
=(1+\varepsilon)(1+\varepsilon^2)(1+\varepsilon^4)\cdots(1+\varepsilon^{2^{r-1}})
=\sum_{\lambda=0}^{2^r-1}\varepsilon^\lambda.
\]
Якщо до $\Pi$ додати ще $\varepsilon^{2^r}=\varepsilon^{-1}=\varepsilon^{p-1}$ (з огляду на $(*)$), то $\Pi+\varepsilon^{p-1}$ стає сумою цілих частин виразу $1+\varepsilon+\varepsilon^2+\cdots+\varepsilon^{p-1}=0$, і тому
\[
        \Pi=-\varepsilon^{-1}.
\]
Тепер і $\varepsilon$ є відносною нормою числа з $P(\varepsilon i)$, а саме числа $\varepsilon^{(p+1)/2}$. Отже, аналогічне твердження справджується для
\[
        \Pi\varepsilon=-1.
\]
Тому $-1$ у $P(\varepsilon)$ подається як сума двох квадратів;\footnote{Для простих чисел форми $8n\pm3$ це міркування вже є в E. Landau, Über die Darstellung definiter Funktionen durch Quadrate, Math. Ann. Bd. 62, с. 272--285.} отже, $P(\varepsilon)$ є полем розщеплення для кватерніонної групи.

Нехай $K$ -- циклічне поле степеня $2^n$, що міститься в $P(\varepsilon)$; ми стверджуємо, що й $K$ є полем розщеплення.\footnote{Вибір циклічного поля розщеплення степеня $2^n$ як підполя поля поділу кола наслідує приклад Гассе.} Нехай $m$ -- індекс кватерніонного тіла щодо $K$ як основного поля; тоді $m=1$ або $m=2$. Але оскільки існує поле розщеплення $P(\varepsilon)$ непарного відносного степеня $(p-1)/2^n$ щодо $K$, другий випадок неможливий; отже, $m=1$. А це означає, що $K$ є полем розщеплення.

Отже, для нескінченно багатьох $n$ існують циклічні поля степеня $2^n$, які є мінімальними полями розщеплення.

\end{document}
